\section{Results}
\label{sec:results}

\subsection{An open, autonomous measurement pipeline}
We instantiate the pipeline of \S\ref{sec:methods} end to end with no human coders and no paid API.
Claim extraction turns each legislative document into a set of empirical claims; for each claim we
retrieve up to five candidate references from OpenAlex \citep{priem2022openalex}, an open index of
scholarly works, so that every reference is a real, dereferenceable record rather than a model-generated
citation; and each (claim, reference) pair is assigned a stance---\textsc{support}, \textsc{refute}, or
\textsc{silent}---by a high-capacity LLM judge (Claude Opus~4.8), a panel of three open-source models,
and two natural-language-inference (NLI) classifiers. All models run locally on a single consumer GPU.
The scientific question---what fraction of legislators' empirical claims the scientific record
supports---reduces to aggregating these stance labels with valid uncertainty, which in turn requires
that the judge be trustworthy. We establish that next.

\subsection{The LLM judge attains human-level accuracy on claim verification}
\label{sec:res-validation}
We validate every judge against expert-human labels on two scientific-claim-verification benchmarks
that bracket the policy setting: SciFact \citep{wadden2020fact}, expert-written claims paired with cited
abstracts (340 pairs), and Climate-FEVER \citep{diggelmann2020climate}, contested real-world claims with
retrieved evidence (class-balanced sample). Both use the identical three-way judgment our pipeline
makes. Each validator predicts a stance for every pair, scored against the human gold
(Table~\ref{tab:validation}, Fig.~\ref{fig:validation}).

On SciFact, the high-capacity judge reaches \textbf{80.3\% accuracy} and \textbf{0.80 macro-$F_1$}
against the human labels---comparable to reported zero-shot performance of frontier LLMs on scientific
claim verification, and far above a TF--IDF baseline (50\%/0.48). Its strongest class is \textsc{refute}
($F_1=0.93$): the judge reliably detects when a study \emph{contradicts} a claim, answering the worry
that an autonomous pipeline cannot surface disconfirming evidence. The two zero-shot NLI classifiers
reach 70\% accuracy ($\textrm{macro-}F_1\approx0.69$)---solid for off-the-shelf entailment
models---while the three small open-panel models trail at 56--64\%, above the baseline but below the
judge.

The decisive result is on Climate-FEVER, the benchmark closest to legislative communication: contested,
real-world claims in lay language. Here the high-capacity judge holds at \textbf{77.8\% accuracy}
($\textrm{macro-}F_1=0.79$), whereas the NLI classifiers collapse to \emph{chance} (50\%, $F_1\approx0.47$).
Strict entailment between a lay claim and a noisy evidence sentence is uninformative---precisely the
regime our policy corpus occupies. The judge thus generalizes to policy-like claims where the cheap
surrogate does not, which is the central reason we anchor estimation on the validated judge and treat
NLI only as a surrogate to be corrected (\S\ref{sec:methods-ppi}). Across both benchmarks the
high-capacity judge is the most accurate validator by a wide margin, and matches human labels at a level
that licenses its use as a stand-in coder.

\begin{table}[t]
\centering
\caption{Validator accuracy and macro-$F_1$ against expert-human gold labels on two
scientific-claim-verification benchmarks. The high-capacity LLM judge matches human labels best;
small open models and zero-shot NLI are close behind; all far exceed the TF--IDF baseline. Because the
benchmark text is human-written, this accuracy cannot be an artifact of self-preference.}
\label{tab:validation}
\begin{tabular}{lcccc}
\toprule
 & \multicolumn{2}{c}{SciFact} & \multicolumn{2}{c}{Climate-FEVER} \\
\cmidrule(lr){2-3}\cmidrule(lr){4-5}
Validator & Accuracy & Macro-$F_1$ & Accuracy & Macro-$F_1$ \\
\midrule
Claude Opus 4.8 (judge)   & \textbf{0.803} & \textbf{0.803} & \textbf{0.778} & \textbf{0.785} \\
Open panel: Qwen2.5-3B     & 0.635 & 0.616 & 0.560 & 0.553 \\
Open panel: Phi-3.5-mini   & 0.559 & 0.486 & 0.627 & 0.617 \\
Open panel: OLMo-2-7B      & 0.600 & 0.478 & 0.522 & 0.492 \\
NLI: DeBERTa-MNLI          & 0.703 & 0.690 & 0.507 & 0.479 \\
NLI: BART-MNLI             & 0.706 & 0.696 & 0.500 & 0.461 \\
TF--IDF (baseline)         & 0.500 & 0.476 & --- & --- \\
\bottomrule
\end{tabular}
\end{table}

\begin{figure}[t]
\centering
\includegraphics[width=0.95\textwidth]{fig_validation.png}
\caption{Accuracy and macro-$F_1$ of each validator against expert-human gold labels, on SciFact
(expert scientific claims) and Climate-FEVER (contested real-world claims). The high-capacity LLM judge
is the most accurate on both; small open models and zero-shot NLI follow; all exceed chance and the
TF--IDF baseline.}
\label{fig:validation}
\end{figure}

\subsection{Self-preference does not explain the judge's accuracy}
\label{sec:res-selfpref}
The benchmark results rule out the central circularity threat. SciFact and Climate-FEVER claims and
evidence are written by humans, so a judge that labels them correctly is assessing the science, not
recognizing model-generated prose: self-preference \citep{panickssery2024llm} cannot manufacture an 80\%
match to human labels on human text (Fig.~\ref{fig:confusion}), and it cannot explain the judge's
accuracy on the policy-adjacent Climate-FEVER claims either. The three independently developed open
models perform the same task well above chance on the same human-written benchmarks (56--64\% on
SciFact, 52--63\% on Climate-FEVER; Table~\ref{tab:validation}), confirming that claim verification is
tractable for models that did not generate the text---the high-capacity judge is simply the most
accurate of them.

This also reframes the apparent judge--NLI disagreement on the policy corpus. There the judge and NLI
agree only weakly ($\kappa=0.15$--$0.21$), which one might read as a problem for the judge. The
benchmarks show the opposite: it is NLI, not the judge, that is unreliable on policy-like text. Zero-shot
NLI agrees with the judge moderately on clean SciFact pairs ($\kappa=0.46$--$0.48$) but is no better than
chance on contested Climate-FEVER claims. The weak policy agreement is therefore a property of strict
entailment under relevance retrieval---an abstract seldom \emph{entails} a lay claim even when it
substantiates it---identified \emph{against human labels} as NLI's failure, not the judge's. This is
exactly why we anchor estimation on the validated judge and treat NLI as a surrogate to be corrected
(\S\ref{sec:methods-ppi}), not as a free-standing measure.

\begin{figure}[t]
\centering
\includegraphics[width=0.5\textwidth]{fig_confusion_scifact.png}
\caption{Confusion matrix of the LLM judge against SciFact human gold. Errors concentrate in the
\textsc{silent}/\textsc{support} boundary; \textsc{refute} is recovered accurately, showing the judge
detects disconfirming evidence.}
\label{fig:confusion}
\end{figure}

\subsection{Application: evidence use in California AI legislation}
\label{sec:res-application}
We apply the validated pipeline to California State Assembly member press releases (2015--2025). Two
findings emerge. \emph{Engagement is large and partisan.} Across the AI-related corpus, Democratic
members engage with AI far more than Republicans on every margin (Fig.~\ref{fig:engage}): they mention
AI in 4.8\% of releases versus 2.6\% for Republicans, and when they do they say more---3.3 AI statements
per document on average versus 1.6. The pipeline retrieves a substantial body of candidate evidence for
the resulting claims.

\emph{Most claims that legislators make are backed by the retrieved science, at similar rates across
parties.} Using the validated judge as the gold anchor and the NLI surrogate over all pairs, the
prediction-powered estimate of the support rate is \textbf{0.71} (95\% CI $[0.57,0.84]$) for Democratic
claims and \textbf{0.54} ($[0.41,0.67]$) for Republican claims (Fig.~\ref{fig:support}). Because the
anchor is a judge we have shown to be \emph{human-level accurate} on this task (\S\ref{sec:res-validation}),
these are validated estimates, not the output of an unchecked model---the key difference from a naive
autonomous measurement. The \emph{level} is genuinely judge-dependent: across the five validators the
Democratic support rate ranges from 0.21 (NLI) to 0.79 (a panel model), which is exactly why benchmark
validation is indispensable---it identifies which judge to trust rather than leaving the analyst to pick
one. What is robust to that choice is the \emph{ordering}: every validator, from NLI through each
open-panel model to the high-capacity judge, rates Democratic claims at least as supported as Republican
ones, so the (modest) partisan gap is not an artifact of any single judge. The partisan intervals
overlap, however: the robust partisan asymmetry is in the \emph{volume} of AI engagement, not in how
well each party's claims are supported.
We read the \emph{level} cautiously: retrieval is relevance-seeking, so it bounds the support rate from
above, and press releases are advocacy, so a high rate reflects selective citation of favorable science.

\begin{figure}[t]
\centering
\includegraphics[width=0.9\textwidth]{fig_ai_engagement.png}
\caption{Democrats engage with AI more often (a) and more intensively (b) than Republicans in
California Assembly press releases.}
\label{fig:engage}
\end{figure}

\subsection{Reproducibility, the retrieval bottleneck, and domain transfer}
\label{sec:res-robustness}
\emph{Reproducible retrieval, and the binding role of retrieval relevance.} Replacing the generative
reference retrieval of earlier pipelines with retrieval from the open OpenAlex index returns only real,
DOI-bearing papers, eliminating the citation-fabrication failure mode of asking a model to name sources.
It also exposes the binding constraint on this kind of measurement. OpenAlex keyword search returns
broad, topically-related papers rather than the specific study a claim invokes, so across \emph{every}
judge---both NLI classifiers and all three open panel models---nearly all pairs are labelled
\textsc{silent} and support rates fall to a few percent (0.6--13\%). The measured support \emph{level} is
thus governed at least as much by how precisely retrieval targets the claim as by the judge: the
0.5--0.7 rates above reflect bespoke, claim-targeted retrieval, whereas broad keyword retrieval surfaces
little claim-specific evidence at all. This is an honest limit of the present pipeline, and it localizes
the remaining work precisely: the \emph{judge} is fully reproducible and human-validated, so it is the
\emph{retrieval} step---a fine-tuned open retriever in place of keyword search or generative
listing---that must be made both reproducible and claim-targeted. The open pipeline is released to
support exactly that. \emph{Domain transfer.} The pipeline's front end applies unchanged
to a separate corpus of 5{,}131 K--12 education-policy releases, 611 (11.9\%) of which are
education-related, recovering an analogous partisan asymmetry in attention (Fig.~\ref{fig:edu}):
Democrats issued 80 releases on social-emotional learning and student mental health versus 1 for
Republicans, and 38 on school safety versus 4. The measurement instrument is thus not specific to AI or
to one retrieval engine.

\begin{figure}[t]
\centering
\includegraphics[width=0.7\textwidth]{fig_support_party.png}
\caption{Claim-support rate by party under each judge (bespoke retrieval). The \emph{level} varies
widely across judges (0.11--0.79 for Democratic claims), but every judge---from zero-shot NLI through
each open-panel model to the validated high-capacity judge (gold, PPI-corrected)---rates Democratic
claims at least as supported as Republican. Benchmark validation identifies which level to trust.}
\label{fig:support}
\end{figure}

\begin{figure}[t]
\centering
\includegraphics[width=0.6\textwidth]{fig_education.png}
\caption{Partisan asymmetry in education-policy attention (counts of press releases), illustrating that
the pipeline's front end generalizes beyond AI.}
\label{fig:edu}
\end{figure}
